% GENERATED FILE. Edit canonical lessons, then run npm run generate:books.
% canonical-source-hash: fnv1a64:c2ad7e030064dc39
% canonical-lessons: JA-C436-hogaraka, JA-C436-akarui, JA-C436-mukuchi, JA-C436-hitomishiri, JA-C436-furumau

\chapter{Cheerful or Quiet}
\label{ch:ja-cheerful-or-quiet}
\hlchaptermodality{436}

\begin{chapteropening}
\textbf{By the end of this chapter:} \emph{I can describe a cheerful or a quiet person.}

Complete the last lesson of chapter 436: I can describe a cheerful or a quiet person.
\end{chapteropening}

\section[\texorpdfstring{\ja{ほがらか} (hogaraka)}{ほがらか (hogaraka)}]{\ja{ほがらか} (hogaraka) — cheerful; sunny-natured}

\label{lesson:JA-C436-hogaraka}



\begin{quote}
Before the new one: say the Japanese for to smile gently, then the Japanese for to fall silent.
\end{quote}

\subsection*{You'll want to know: \ja{ほがらか}}
\textbf{\ja{ほがらか}} — \emph{hogaraka} — \textquotedblleft{}cheerful; sunny-natured\textquotedblright{}.

\begin{quote}
\textbf{\ja{あねは} \ja{いつも} \ja{ほがらかです。}} — \emph{ane wa itsumo hogaraka desu} — \textquotedblleft{}My elder sister is cheerful all the time.\textquotedblright{}
\end{quote}

A na-adjective: \ja{な} before a noun.

\begin{cousinweb}
\textbf{\ja{ほがらか}} (\emph{hogaraka}) is a native Japanese word (wago), inherited rather than borrowed from Chinese; its kanji spelling was attached to it later (JMdict).
\end{cousinweb}

\begin{grammarlens}[title={What you've built}]
One more word for describing what you lived and felt.
\end{grammarlens}

\subsection*{Your turn}
\begin{itemize}
\raggedright
  \item \emph{Say it:} \emph{hogaraka}
  \item \emph{Say it:} \emph{hogaraka}, once more
  \item \emph{Say it:} \emph{ane wa itsumo hogaraka desu} — My elder sister is cheerful all the time.
  \item \emph{Recall:} say the Japanese for to smile gently, then the Japanese for to fall silent, then say \emph{hogaraka} again
\end{itemize}

\begin{tcolorbox}[breakable,colback=teal!4,colframe=teal!35!black,title={Before you move on}]
What does \ja{ほがらか} mean? (\textquotedblleft{}Cheerful; sunny-natured\textquotedblright{}.) Say it once more.
\end{tcolorbox}
\section[\texorpdfstring{\ja{あかるい} (akarui)}{あかるい (akarui)}]{\ja{あかるい} (akarui) — bright; light}

\label{lesson:JA-C436-akarui}



\begin{quote}
Before the new one: say the Japanese for to fall silent, then the Japanese for cheerful.
\end{quote}

\subsection*{You'll want to know: \ja{あかるい}}
\textbf{\ja{あかるい}} — \emph{akarui} — \textquotedblleft{}bright; light\textquotedblright{}.

\begin{quote}
\textbf{\ja{かれは} \ja{あかるい} \ja{せいかくです。}} — \emph{kare wa akarui seikaku desu} — \textquotedblleft{}He has a bright personality.\textquotedblright{}
\end{quote}

An i-adjective.

\begin{cousinweb}
\textbf{\ja{あかるい}} (\emph{akarui}) is a native Japanese word (wago), inherited rather than borrowed from Chinese; its kanji spelling was attached to it later (JMdict).
\end{cousinweb}

\begin{grammarlens}[title={What you've built}]
One more word for describing what you lived and felt.
\end{grammarlens}

\subsection*{Your turn}
\begin{itemize}
\raggedright
  \item \emph{Say it:} \emph{akarui}
  \item \emph{Say it:} \emph{akarui}, once more
  \item \emph{Say it:} \emph{kare wa akarui seikaku desu} — He has a bright personality.
  \item \emph{Recall:} say the Japanese for to fall silent, then the Japanese for cheerful, then say \emph{akarui} again
\end{itemize}

\begin{tcolorbox}[breakable,colback=teal!4,colframe=teal!35!black,title={Before you move on}]
What does \ja{あかるい} mean? (\textquotedblleft{}Bright; light\textquotedblright{}.) Say it once more.
\end{tcolorbox}
\section[\texorpdfstring{\ja{むくち} (mukuchi)}{むくち (mukuchi)}]{\ja{むくち} (mukuchi) — taciturn; of few words}

\label{lesson:JA-C436-mukuchi}



\begin{quote}
Before the new one: say the Japanese for cheerful, then the Japanese for bright.
\end{quote}

\subsection*{You'll want to know: \ja{むくち}}
\textbf{\ja{むくち}} — \emph{mukuchi} — \textquotedblleft{}taciturn; of few words\textquotedblright{}.

\begin{quote}
\textbf{\ja{ちちは} \ja{あさは} \ja{むくちです。}} — \emph{chichi wa asa wa mukuchi desu} — \textquotedblleft{}My father says little in the mornings.\textquotedblright{}
\end{quote}

A na-adjective: \ja{な} before a noun.

\begin{cousinweb}
\textbf{\ja{むくち}} (\emph{mukuchi}) is a mixed word: Sino-Japanese \emph{mu}, \textquotedblleft{}nothing\textquotedblright{} (Middle Chinese \emph{mju}), borrowed from Chinese, joined to native \emph{kuchi}, \textquotedblleft{}mouth\textquotedblright{}, the kun'yomi of its second kanji (KANJIDIC2; Middle Chinese in Baxter's transcription).
\end{cousinweb}

\begin{grammarlens}[title={What you've built}]
One more word for describing what you lived and felt.
\end{grammarlens}

\subsection*{Your turn}
\begin{itemize}
\raggedright
  \item \emph{Say it:} \emph{mukuchi}
  \item \emph{Say it:} \emph{mukuchi}, once more
  \item \emph{Say it:} \emph{chichi wa asa wa mukuchi desu} — My father says little in the mornings.
  \item \emph{Recall:} say the Japanese for cheerful, then the Japanese for bright, then say \emph{mukuchi} again
\end{itemize}

\begin{tcolorbox}[breakable,colback=teal!4,colframe=teal!35!black,title={Before you move on}]
What does \ja{むくち} mean? (\textquotedblleft{}Taciturn; of few words\textquotedblright{}.) Say it once more.
\end{tcolorbox}
\section[\texorpdfstring{\ja{ひとみしり} (hitomishiri)}{ひとみしり (hitomishiri)}]{\ja{ひとみしり} (hitomishiri) — shy with strangers}

\label{lesson:JA-C436-hitomishiri}



\begin{quote}
Before the new one: say the Japanese for bright, then the Japanese for taciturn.
\end{quote}

\subsection*{You'll want to know: \ja{ひとみしり}}
\textbf{\ja{ひとみしり}} — \emph{hitomishiri} — \textquotedblleft{}shy with strangers\textquotedblright{}.

\begin{quote}
\textbf{\ja{むすめは} \ja{すこし} \ja{ひとみしりです。}} — \emph{musume wa sukoshi hitomishiri desu} — \textquotedblleft{}My daughter is a little shy with strangers.\textquotedblright{}
\end{quote}

A na-adjective: \ja{な} before a noun.

\begin{cousinweb}
\textbf{\ja{ひとみしり}} (\emph{hitomishiri}) is a native Japanese compound (wago), inherited rather than borrowed from Chinese: \emph{hito}, \textquotedblleft{}person\textquotedblright{}, \emph{mi}, \textquotedblleft{}seeing (miru)\textquotedblright{}, and \emph{shiri}, \textquotedblleft{}knowing (shiru)\textquotedblright{} (JMdict, KANJIDIC2).
\end{cousinweb}

\begin{grammarlens}[title={What you've built}]
One more word for describing what you lived and felt.
\end{grammarlens}

\subsection*{Your turn}
\begin{itemize}
\raggedright
  \item \emph{Say it:} \emph{hitomishiri}
  \item \emph{Say it:} \emph{hitomishiri}, once more
  \item \emph{Say it:} \emph{musume wa sukoshi hitomishiri desu} — My daughter is a little shy with strangers.
  \item \emph{Recall:} say the Japanese for bright, then the Japanese for taciturn, then say \emph{hitomishiri} again
\end{itemize}

\begin{tcolorbox}[breakable,colback=teal!4,colframe=teal!35!black,title={Before you move on}]
What does \ja{ひとみしり} mean? (\textquotedblleft{}Shy with strangers\textquotedblright{}.) Say it once more.
\end{tcolorbox}
\section[\texorpdfstring{\ja{ふるまう} (furumau)}{ふるまう (furumau)}]{\ja{ふるまう} (furumau) — to behave; to conduct oneself}

\label{lesson:JA-C436-furumau}



\begin{quote}
Before the new one: say the Japanese for taciturn, then the Japanese for shy with strangers.
\end{quote}

\subsection*{You'll want to know: \ja{ふるまう}}
\textbf{\ja{ふるまう}} — \emph{furumau} — \textquotedblleft{}to behave; to conduct oneself\textquotedblright{}.

\begin{quote}
\textbf{\ja{むすこは} \ja{あかるく} \ja{ふるまいました。}} — \emph{musuko wa akaruku furumaimashita} — \textquotedblleft{}My son acted cheerful.\textquotedblright{}
\end{quote}

Polite \textbf{\ja{ふるまいます}}, \emph{furumaimasu}; a godan verb. Past (polite): \textbf{\ja{こどもが} \ja{おとなしく} \ja{ふるまいました。}}, \emph{kodomo ga otonashiku furumaimashita}.

\begin{cousinweb}
\textbf{\ja{ふるまう}} (\emph{furumau}) is a native Japanese word (wago), inherited rather than borrowed from Chinese; its kanji spelling was attached to it later (JMdict).
\end{cousinweb}

\begin{grammarlens}[title={What you've built}]
One more word for describing what you lived and felt.
\end{grammarlens}

\subsection*{Your turn}
\begin{itemize}
\raggedright
  \item \emph{Say it:} \emph{furumau}
  \item \emph{Say it:} \emph{furumau}, once more
  \item \emph{Say it:} \emph{musuko wa akaruku furumaimashita} — My son acted cheerful.
  \item \emph{Recall:} say the Japanese for taciturn, then the Japanese for shy with strangers, then say \emph{furumau} again
\end{itemize}

\begin{tcolorbox}[breakable,colback=teal!4,colframe=teal!35!black,title={Before you move on}]
What does \ja{ふるまう} mean? (\textquotedblleft{}To behave; to conduct oneself\textquotedblright{}.) Say it once more.
\end{tcolorbox}
